% Purpose: Academic introduction to earthquake mechanics, wave
% propagation physics, regional seismotectonics, SMINO network,
% classical vs. ML picking, and research roadmap.
% Status: REVISED
% Source-of-truth: OGS/src/, SMINO documentation, literature
% citations in lib.bib.

\label{chap:introduction}

\section{Seismological Foundations and Wave Propagation Physics}
\label{sec:wave_physics}

Earthquakes represent transient mechanical disturbances generated by the sudden
release of accumulated elastic strain energy within the Earth's lithosphere.
The lithosphere behaves as a brittle-elastic medium under tectonic differential
stress; when shear stresses acting along a pre-existing fault plane exceed the
frictional resistance governed by Byerlee's law or rate-and-state friction,
dynamic shear failure nucleates and ruptures along the fault surface
\parencite{Byerlee1978,Dieterich1979}. This catastrophic dislocation radiates
transient elastodynamic stress waves that propagate through the heterogeneous
interior and along the free surface of the Earth
\parencite{KanamoriBrodsky2004}.

In a continuous, isotropic, linearly elastic medium characterized by mass
density $\rho(\mathbf{x})$ and Lamé elastic moduli $\lambda(\mathbf{x})$ and
$\mu(\mathbf{x})$, the infinitesimal elastodynamic equation of motion
expressing conservation of linear momentum is formulated as:
\begin{equation}
  \rho(\mathbf{x}) \frac{\partial^2 \mathbf{u}}{\partial t^2}(\mathbf{x}, t) =
  \nabla \cdot \bm{\sigma}(\mathbf{x}, t) + \mathbf{f}(\mathbf{x}, t),
  \label{eq:momentum_conservation}
\end{equation}
where $\mathbf{u}(\mathbf{x}, t) = [u_1, u_2, u_3]^\top$ denotes the
displacement vector field, $\bm{\sigma}(\mathbf{x}, t)$ is the symmetric
Cauchy stress tensor, and $\mathbf{f}(\mathbf{x}, t)$ represents external body
force distributions per unit volume (such as equivalent double-couple source
representations of fault rupture).

Under the generalized Hooke's law for isotropic linear elasticity, the
constitutive stress-strain relation binds Cauchy stress to the
infinitesimal strain tensor $\bm{\varepsilon}$:
\begin{equation}
  \bm{\sigma} \coloneqq \lambda (\nabla \cdot \mathbf{u}) \mathbf{I} +
  2\mu \bm{\varepsilon} \equiv \lambda (\operatorname{div}
  \mathbf{u}) \mathbf{I} +
  \mu \left( \nabla \mathbf{u} + (\nabla \mathbf{u})^\top \right),
  \label{eq:hooke_law_tensor}
\end{equation}
or in component index notation:
\begin{equation}
  \sigma_{ij} \coloneqq \lambda \delta_{ij} \epsilon_{kk} + 2\mu \epsilon_{ij},
  \label{eq:hooke_law}
\end{equation}
where $\delta_{ij}$ is the Kronecker delta, $\epsilon_{ij} \coloneqq
\frac{1}{2}\left(\partial_{x_j} u_i + \partial_{x_i} u_j\right)$ denotes the
symmetric infinitesimal strain tensor, and Einstein summation convention over
repeated indices is implied.

For the \textbf{homogeneous-medium} specialization used below, let $\rho > 0$
and $\lambda,\mu > 0$ be spatially constant ($\nabla \rho = \nabla \lambda =
\nabla \mu = \mathbf{0}$). Substituting the constitutive law in
\cref{eq:hooke_law} into the equation of motion in
\cref{eq:momentum_conservation} yields the Navier--Cauchy elastodynamic wave
equation:
\begin{equation}
  \rho \ddot{\mathbf{u}} = (\lambda +
  \mu)\nabla(\nabla \cdot \mathbf{u}) + \mu \nabla^2 \mathbf{u} + \mathbf{f},
  \label{eq:navier_cauchy}
\end{equation}

For \textbf{source-free} body-wave modes, set $\mathbf{f} = \mathbf{0}$. On an
\textbf{unbounded domain}, \textbf{assume} that the displacement and potentials
decay sufficiently rapidly at infinity. With the Coulomb gauge
($\nabla \cdot \bm{\psi} \equiv 0$), Helmholtz's decomposition expresses the
vector displacement field as the sum of an irrotational (curl-free) scalar
potential $\phi(\mathbf{x}, t)$ and a solenoidal (divergence-free) vector
potential $\bm{\psi}(\mathbf{x}, t)$:
\begin{equation}
  \mathbf{u} \coloneqq \nabla \phi + \nabla \times \bm{\psi}, \quad
  \text{with} \quad \nabla \cdot \bm{\psi} \equiv 0.
  \label{eq:helmholtz}
\end{equation}
The Coulomb gauge condition ($\nabla \cdot \bm{\psi} \equiv 0$) and decay
conditions, remove harmonic residuals. Applying the divergence operator
and curl operator to \cref{eq:helmholtz}, and exploiting the fundamental vector
nullity identities $\nabla \times (\nabla \phi) \equiv \mathbf{0}$ and
$\nabla \cdot (\nabla \times \bm{\psi}) \equiv 0$, decomposes the wavefield
into orthogonal kinematic projections:
\begin{equation}
  \nabla \cdot \mathbf{u} = \nabla^2 \phi, \quad
  \nabla \times \mathbf{u} = -\nabla^2 \bm{\psi}.
  \label{eq:helmholtz_projections}
\end{equation}
Substituting \cref{eq:helmholtz,eq:helmholtz_projections} into the source-free
Navier--Cauchy equation in \cref{eq:navier_cauchy} decouples the elastodynamic
wavefield into two independent wave equations governing the temporal kinematic
acceleration rates of the potentials:
\begin{equation}
  \ddot{\phi} = \alpha^2 \nabla^2 \phi, \quad
  \ddot{\bm{\psi}} = \beta^2 \nabla^2 \bm{\psi},
  \label{eq:decoupled_wave_equations}
\end{equation}
where $\ddot{\phi} \equiv \frac{\partial^2 \phi}{\partial t^2}$ and
$\ddot{\bm{\psi}} \equiv \frac{\partial^2 \bm{\psi}}{\partial t^2}$ denote the
second-order temporal kinematic rates of change for the compressional and
equivoluminal potential fields, and $\alpha$ and $\beta$ represent the
longitudinal compressional (Primary, P) and transverse shear (Secondary, S)
wave propagation phase velocities, respectively:
\begin{equation}
  \alpha \coloneqq \sqrt{\frac{\lambda + 2\mu}{\rho}}, \quad \beta \coloneqq
  \sqrt{\frac{\mu}{\rho}}.
  \label{eq:wave_velocities}
\end{equation}

The physical characteristics and particle motion trajectories of these seismic
phases govern the observational records captured by modern seismological sensor
systems:
\begin{enumerate}[label=(\alph*)]
  \item \textbf{Primary (P) Waves}: Compressional or dilatational longitudinal
    waves that propagate via volumetric deformation ($\nabla \cdot \mathbf{u}
    \not\equiv 0, \; \nabla \times \mathbf{u} \equiv \mathbf{0}$). Particle
    motion is strictly parallel to the direction of wave propagation. Because
    $\lambda > 0$ and $\mu > 0$; $\alpha > \beta$ holds universally, ensuring
    that P-waves represent the fastest seismic phase and form the first onset
    arrival observed on seismograms.
  \item \textbf{Secondary (S) Waves}: Equivoluminal transverse shear waves that
    propagate via shear deformation without volume change
    ($\nabla \cdot \mathbf{u} \equiv 0, \; \nabla \times \mathbf{u} \not\equiv
    \mathbf{0}$). Particle displacement occurs perpendicular to the ray
    propagation vector. In a Cartesian coordinate system aligned with the
    incidence plane, S-waves are decomposed into \ac{SV} and \ac{SH}.
    Because fluids have zero shear modulus ($\mu = 0$), S-waves propagate
    exclusively through solid elastic media.
\end{enumerate}

In an ideal Poissonian solid ($\lambda = \mu$, equivalent to Poisson's ratio
$\nu = \num{0.25}$), the kinematic velocity ratio satisfies:
\begin{equation}
  \frac{\alpha}{\beta} \equiv \sqrt{\frac{\lambda + 2\mu}{\mu}} = \sqrt{3}
  \approx \num{1.732}.
  \label{eq:poisson_ratio}
\end{equation}
In complex geological crustal formations, lithological variations, microcrack
porosity, and fluid saturation cause the observed
$V_{\mathrm{P}}/V_{\mathrm{S}}$ ratio to vary typically between $\num{1.65}$
and $\num{1.90}$ \placeholder{CITE}.

At \textbf{discontinuous} geological interfaces and the Earth's
\textbf{stress-free} surface ($\sigma_{i3}|_{z = 0} = 0$), mode conversion and
interference between body waves generate \textbf{surface waves}. These
propagate horizontally along the crustal waveguide with amplitudes decaying
exponentially with depth:
\begin{itemize}
  \item \textbf{Rayleigh Waves}: Resulting from constructive P-\ac{SV}
    interference at the free boundary, exhibiting retrograde elliptical
    particle motion in the vertical-radial propagation plane, with phase
    velocities $c_{\mathrm{R}} \approx \num{0.92}\beta$.
  \item \textbf{Love Waves}: Resulting from constructive \ac{SH} wave trapping
    within low-velocity surface layers overlying a higher-velocity half-space,
    exhibiting purely horizontal transverse motion perpendicular to the
    propagation path.
\end{itemize}

Modern digital seismological stations record these vector ground motions along
three mutually orthogonal spatial components: Vertical (Z), North-South (N),
and East-West (E). Typically, modern seismic stations of local networks deploy
high-gain broadband seismometers to measure ground velocity
$\dot{\mathbf{u}}(t)$ over flat bandwidths (\qtyrange{50}{100}{\hertz})
and strong-motion accelerometers to measure ground acceleration
$\ddot{\mathbf{u}}(t)$. Data are continuously digitized at high sampling rates
(generally \qtyrange{100}{200}{\hertz}) and distributed in standard miniSEED
(MSEED) format, with instrumental response information and geodetic station
coordinates formally cataloged in StationXML schemas.

\section{Seismotectonic Context of Northeastern Italy and the 1976 Legacy}
\label{sec:tectonics}

Northeastern Italy, embracing the \acf{FVG} and Veneto regions, represents one
of the most tectonically active and seismically hazardous sectors of the entire
European Alpine arc. The geodynamic evolution of the region is governed by the
active collision and northward indentation of the Adriatic microplate (Adria)
into the stable Eurasian plate. Geodetic measurements from continuous GNSS
networks indicate that the Adria microplate rotates counterclockwise around an
Eulerian pole in the Western Po Plain, producing active crustal shortening
across the Eastern Southern Alps at rates of
\qtyrange{1.5}{3.0}{\milli\meter\per\year} in a NNW-SSE orientation
\parencite{AndersonJackson1987,Handy2014,Tunini2024}.

Structurally, Northeastern Italy constitutes a complex tectonic junction where
two major collisional orogenic belts converge:
\begin{enumerate}
  \item The \textbf{Eastern Southern Alps}: A south-verging thrust-and-fold
    retro-wedge characterized by ENE-WSW to WNW-ESE striking thrust faults and
    ramp-anticlines (including the Periadriatic Lineament, the Fella-Sava
    fault, the Maniago thrust, and the Buia-Tricesimo thrust system) that
    accommodate crustal shortening through north-dipping blind and exposed
    thrusting \parencite{Bressan2018}.
  \item The \textbf{External Dinarides}: A dextral strike-slip structural
    system extending from Slovenia into Eastern Friuli, dominated by NW-SE
    trending transcurrent fault structures such as the Idrija and Ravne fault
    networks \parencite{Bressan1998}.
\end{enumerate}

This tectonic configuration results in concentrated seismicity characterized by
shallow hypocenters (typically \qtyrange{5}{15}{\kilo\meter} depth),
predominantly compressive thrust-faulting mechanisms with subordinate
strike-slip kinematics, and destructive earthquake sequences
\parencite{Sarao2021}.

The strongest seismic event recorded in this region is the catastrophic
\textbf{1976 Friuli earthquake sequence}. On May 6, 1976, at 20:00:12 UTC, a
major earthquake of moment magnitude $M_{\mathrm{w}} = \num{6.5}$ nucleated
near Gemona del Friuli, producing epicentral intensities of $I_0 =
\mathrm{IX}\text{--}\mathrm{X}$ on the \ac{MCS} scale \parencite{Locati2022}.
The mainshock ruptured blind thrust faults along the Southalpine front
\parencite{Aoudia2000}. The sequence was compounded in September 1976 by two
devastating aftershocks on September 15 ($M_{\mathrm{w}} = 5.6$ and
$M_{\mathrm{w}} = \num{6.0}$). In total, the 1976 Friuli seismic crisis claimed
\num{989} lives, injured over \num{2800} people, left over \num{100000}
homeless, and destroyed entire medieval villages across the Tagliamento and
Fella river basins.

The tragedy of the 1976 disaster catalyzed a profound transformation in
European seismology, earthquake engineering, civil protection organization, and
disaster response legislation. Crucially, it led the Regional Government of
\ac{FVG} to mandate the establishment of a dedicated permanent seismological
monitoring infrastructure. In 1977, the \acf{OGS} established the \acf{CRS} in
Udine. Since its inception, \ac{OGS} \ac{CRS} has maintained continuous,
high-resolution seismic monitoring across Northeastern Italy, serving as the
official technical-scientific body supporting the regional Civil Protection
agencies of \ac{FVG}.

\section{The \ac{OGS}-\acs{SMINO} Monitoring Network Infrastructure}
\label{sec:smino_network}

Over five decades of technological development, the initial sparse analog
network deployed after the 1976 earthquake evolved into the modern \acf{SMINO}.
The \ac{SMINO} monitoring infrastructure represents a state-of-the-art,
multi-parametric, real-time observation system covering \ac{FVG}, Veneto, and
surrounding border areas.

The regional monitoring infrastructure integrates:
\begin{itemize}
  \item A dense permanent backbone of digital three component broadband and
    very broadband seismometers.
  \item High-dynamic-range strong-motion accelerometers collocated at broadband
    sites and deployed across urban centers and critical engineering
    infrastructures to ensure on-scale recordings of strong ground motions.
  \item Cross-border data exchange with neighboring Italian, Austrian, and
    Slovenian institutions, creating a unified transboundary seismic monitoring
    domain.
\end{itemize}

As described in \cref{chap:datasets}, the complete regional monitoring aperture
analyzed in this study encompasses \placeholder{NUM STATIONS} stations
originating from \placeholder{NUM NETWORKS} distinct regional,
national, and international networks (\cref{tab:network_inventory} provides a
  systematic synthesis of these network assets, followed by formal
bibliographic and data repository citations). The continuous
recording from these seismic stations provides the empirical
foundation for our \ac{ML} and computational investigations.

\section{Classical Automated Picking}
\label{sec:classical_picking}

Transforming continuous, multi-channel seismic streams into discrete, verified
earthquake catalogs requires identifying the exact onset times of P- and S-wave
phase arrivals (phase picking) and grouping arrivals originating from the same
physical hypocentral rupture (phase association).

For over four decades, automated seismological processing has relied on
classical heuristic algorithms:
\begin{enumerate}
  \item \textbf{\acf{STA/LTA}}: The method calculates the ratio of energy or
    amplitude within a short sliding time window ($T_{\mathrm{STA}}$) to that
    within a preceding long time window ($T_{\mathrm{LTA}}$) over a
    characteristic function $\mathrm{CF}(t)$:
    \begin{equation}
      R(t) \coloneqq \frac{\mathrm{STA}(t)}{\mathrm{LTA}(t)} =
      \frac{\frac{1}{T_{\mathrm{STA}}}\int_{t - T_{\mathrm{STA}}}^t
      \mathrm{CF}(\tau)\,\mathrm{d}\tau}{\frac{1}{T_{\mathrm{LTA}}}\int_{t
      - T_{\mathrm{LTA}}}^t\mathrm{CF}(\tau)\,\mathrm{d}\tau},
      \label{eq:stalta}
    \end{equation}
    where $\mathrm{CF}(t)$ is commonly defined as $y(t)^2$ or $y(t)^2 +
    C \dot{y}(t)^2$. A pick trigger is declared when $R(t) \ge
    \eta_{\mathrm{on}}$, and deactivated when $R(t) \le \eta_{\mathrm{off}}$.
  \item \textbf{Autoregressive \ac{AIC}}: Applied to seismic phase timing by
    \textcite{LEONARD1999247}, the \ac{AIC} models the seismogram before and
    after an arbitrary split point $k$ as two distinct stationary
    autoregressive processes. The minimum of the \ac{AIC} function identifies
    the optimal transition between background noise and seismic onset:
    \begin{equation}
      \mathrm{AIC}(k) \coloneqq k \log\bigl(\hat{\sigma}_1^2\bigr) + (N - k)
      \log\bigl(\hat{\sigma}_2^2\bigr),
      \label{eq:aic}
    \end{equation}
    where $\hat{\sigma}_1^2$ and $\hat{\sigma}_2^2$ are the variance estimators
    of the prediction errors for intervals $[1, k]$ and $[k + 1, N]$,
    respectively.
\end{enumerate}

Despite their enduring historical utility, classical algorithms suffer from
fundamental physical and algorithmic limitations:
\begin{itemize}
  \item \textbf{Vulnerability to Noise and False Triggers}: \ac{STA/LTA} is
    highly sensitive to non-stationary transient noise, cultural vibrations,
    electrical glitches, and meteorological disturbances, causing false
    triggers in densely populated or industrial valleys.
  \item \textbf{Emergent and Low-SNR Degradation}: Onsets of small-magnitude
    earthquakes ($M_{\mathrm{L}} < \num{2.0}$) or emergent arrivals traveling
    through attenuative, heterogeneous crustal paths exhibit low \acp{SNR}.
  \item \textbf{Phase Misidentification}: Neither \ac{STA/LTA} nor \ac{AIC}
    natively distinguishes between compressional (P) and shear (S)
    phases without ad-hoc, brittle horizontal-to-vertical amplitude
    ratios or polarization filters.
  \item \textbf{Combinatorial Association Bottlenecks}: Grid-search associators
    that cluster picks by evaluating travel-time residuals across spatial
    velocity grids exhibit exponential computational scaling
    $\mathcal{O}(N_{\mathrm{picks}}^k)$ with pick volume. When spurious noise
    picks enter the pipeline, association algorithms suffer combinatoric
    explosion or form hallucinated, unphysical hypocenters.
\end{itemize}

Because of these limitations, operational seismological observatories have
historically required dedicated human analysts to inspect, refine, or manually
pick phase arrivals for every candidate event. This manual intervention
introduces severe operational latency (often hours to days), creates subjective
timing inconsistencies between analysts.

\section{The Big Data Era and Deep Learning in Seismology}
\label{sec:dl_seismology}

Over the past decade, the rapid expansion of permanent seismic networks, dense
nodal geophone arrays, and real-time transboundary telemetry has propelled
observational seismology into the era of ``Big Data''. Continuous seismic
archives at regional and national data centers ingest tens of gigabytes per day
and routinely accumulate terabytes to petabytes of continuous waveform data
annually. Analyzing these continuous, multi-channel data streams exhaustively
is beyond human manual processing capacity.

This data collection has coincided with revolutionary advances in \acf{DL}.
Deep neural networks possess the unique ability to learn rich, non-linear,
hierarchical feature representations directly from continuous sensor time
series without requiring manual feature engineering. Applied to seismic
waveforms, deep convolutional and attention-based architectures learn the
multi-scale spatio-temporal signatures of wave propagation, dispersion, and
noise textures \parencite{MousaviBeroza2022}.

Key architectural breakthroughs include among others:
\begin{itemize}
  \item \textbf{\acl{PhaseNet}}: A deep 1D \ac{CNN} with a U-Net architecture,
    designed for precise P- and S-wave arrival time picking on three component
    seismograms.
  \item \textbf{\acl{EQT}}: A hierarchical \ac{DL} architecture combining
    \ac{CNN} layers, bidirectional \ac{LSTM}, and multi-head self-attention
    mechanisms to capture local waveform patterns and long-range temporal
    dependencies across \qty{60}{\second} windows, simultaneously outputting
    detection probabilities for earthquake presence, P-wave arrival, and S-wave
    arrival.
\end{itemize}

To overcome software fragmentation and provide standardized interfaces for
benchmarking, \textcite{SeisBench} introduced \textbf{SeisBench}, an
open-source platform offering unified Python \acp{API} for \ac{ML} in
seismology. SeisBench provides standardized data loaders, benchmark training
datasets, and pre-trained weights for state-of-the-art architectures,
dramatically enhancing the reproducibility of \ac{DL} seismic workflows. In
this thesis, SeisBench serves as the core foundation for evaluating pre-trained
phase pickers across diverse training datasets.

\section{The March 27, 2024 Socchieve Earthquake Sequence}
\label{sec:socchieve_sequence}

To test the operational efficacy, robustness, and scientific fidelity of
\ac{DL} pipelines under demanding real-world conditions, this research focuses
on an intense seismic sequence that occurred in Friuli in early 2024.

On \textbf{March 27, 2024, at 21:19:37 UTC} (22:19:37 local time), a magnitude
$M_{\mathrm{L}} = \num{4.6 \pm 0.3}$ ($M_{\mathrm{w}} = \num{4.5}$) earthquake
occurred in the Carnic Prealps. The earthquake was strongly felt across the
entire \ac{FVG} and Veneto regions, as well as in neighboring Austria and
Slovenia.

This earthquake constitutes one of the \textbf{most energetic tectonic event to
occur in Friuli since the 1976--1980 crisis}. It initiated a vigorous
aftershock sequence comprising hundreds of minor earthquakes and microseismic
events that decayed over subsequent weeks and months. The spatial proximity of
dense seismic networks, renders the Socchieve sequence a natural testbed for
stress-testing \ac{ML} detection, phase association, and hypocentral relocation
against an expert analyst ground truth.

\section{Research Objectives and Scope of the Dissertation}
\label{sec:objectives}

The scientific objective of this doctoral dissertation is to construct,
optimize, deploy, and rigorously validate a complete, automated, procedural
\ac{ML} pipeline for earthquake pick detection, phase association, and event
characterization, specifically engineered for scalable execution on \ac{HPC}
infrastructures.

The key research questions addressed in this thesis are:
\begin{enumerate}
  \item \textbf{Domain Adaptation and Generalization}: How effectively do
    \ac{DL} phase pickers (\ac{PhaseNet} and \ac{EQT}) trained on large
    external datasets (\ac{INSTANCE} in Italy, \ac{STEAD} globally, \ac{SCEDC}
    in California) generalize to the complex geological and noise environment
    of Northeastern Italy?
  \item \textbf{Associator Scaling and Performance}: How do advanced phase
    association algorithms (specifically the Bayesian \acf{GaMMA} and the
    \acf{PyOcto}) compare in association recall, computational efficiency, and
    robustness against \ac{FP} pick contamination?
  \item \textbf{Physically Bounded Performance Validation}: How can
    model-generated earthquake catalogs be compared against manual analyst
    reference catalogs without introducing spatiotemporal bias or duplicate
    assignment artifacts?
  \item \textbf{HPC Scalability and Operational Viability}: Can the entire
    processing workflow \ie{from continuous multi-client \ac{FDSN} waveform
    ingestion to deep inference, association, and relocation} be parallelized
    across modern \ac{GPU}-accelerated supercomputing architectures (such as
    the \ac{CINECA} Leonardo Tier-0 system) to support near-real-time
    operational monitoring and rapid historical archive reprocessing?
\end{enumerate}

\section{Dissertation Structure and Chapter Roadmap}
\label{sec:roadmap}

The remainder of this dissertation is organized into eight chapters followed by
a concluding synthesis, a dedicated Data and Resources statement, and technical
appendices:
\begin{itemize}
  \item \textbf{\Cref{chap:datasets} (Datasets)}: Details the operational case
    study data acquired across Northeastern Italy evaluating over twenty years
    of observational data (\textbf{January 1st, 2005 to December 31st, 2025})
    across \placeholder{NUM STATIONS} stations, (\placeholder{NUM EVENTS}
    relocated \ac{OGS} reference events with \placeholder{NUM PICKS} picks) and
    reviews the four benchmark training datasets (\ac{INSTANCE}, \ac{STEAD},
    \ac{SCEDC}, and Original sets).
  \item \textbf{\Cref{chap:data_engineering} (Data Engineering)}: Presents the
    software architecture of the pipeline, detailing the central
    \texttt{ogsconstants} configuration, weight-to-uncertainty mappings,
    automated \ac{FDSN} multi-client waveform downloading, multi-format legacy
    catalog parsing (\texttt{.hpl}, \texttt{.dat}, \texttt{.txt}, and
    \texttt{.pun}), catalog aggregation, geodetic math, and visualization
    frameworks.
  \item \textbf{\Cref{chap:computational_pipeline} (Computational Pipeline)}:
    Formulates the complete procedural workflow, detailing continuous waveform
    preprocessing, \ac{DL} phase picking inference, Bayesian (\ac{GaMMA})
    and octree (\ac{PyOcto}) phase association, 1D \ac{NonLinLoc} probabilistic
    relocation, and local magnitude ($M_{\mathrm{L}}$) estimation.
  \item \textbf{\Cref{chap:pipeline_evaluation} (Pipeline Evaluation
    Framework)}: Develops the formal evaluation methodology, defining the
    \acf{BPGMA} for both picks and events, the construction of domain-specific
    confusion matrices, and the mathematical derivation of recall, precision,
    and \acp{F1}.
  \item \textbf{\Cref{chap:hpc} (HPC Implementation)}: Analyzes the \ac{HPC}
    implementation deployed on the \ac{CINECA} Leonardo supercomputer (EuroHPC
    Booster partition), detailing SLURM job orchestration, \ac{MPI} and
    \ac{GPU} process allocation, parallel I/O throughput, and execution scaling
    profiles.
  \item \textbf{\Cref{chap:training} (Training)}: Details the \ac{DL} training
    pipeline, loss functions, optimization, and \ac{HPC} scaling for model
    fine-tuning and regional domain adaptation on continuous waveform archives.
  \item \textbf{\Cref{chap:clustering} (Seismic Swarm and Density-Based
    Clustering)}: Develops density-based and manifold-learning clustering
    formulations, information-theoretic cluster validity indices, and
    epidemic-type aftershock sequence modeling for resolving fine
    spatiotemporal seismicity structures and swarm dynamics.
  \item \textbf{\Cref{chap:results} (Results)}: Delivers a comprehensive
    empirical evaluation of the benchmarking campaign on the Northeastern Italy
    dataset, presenting pick detection recall curves, timing residual
    distributions, hypocentral relocation offsets, magnitude fidelity, and
    comparative associator analysis.
  \item \textbf{\Cref{chap:conclusion} (Conclusion)}: Synthesizes the core
    scientific and engineering findings, discusses the implications for modern
    operational seismological observatories, outlines current limitations, and
    charts future directions for \ac{AI}-driven seismological monitoring.
  \item \textbf{Data and Resources}: Declares the persistent \acp{DOI},
    software repositories, and access endpoints for all code, models, and
    catalogs developed in this work.
  \item \textbf{Appendix}: Documents supplementary mathematical derivations,
    configuration schemas, and detailed per-station performance tables.
\end{itemize}
